\documentclass[11pt, a4paper]{article}
\usepackage[T1]{fontenc}
\usepackage[utf8]{inputenc}
\usepackage[a4paper, top=2.5cm, bottom=2.5cm, left=2.5cm, right=2.5cm]{geometry}
\usepackage{booktabs}
\usepackage{tabularx}
\usepackage{enumitem}
\usepackage[hidelinks]{hyperref}
\usepackage{parskip}
\usepackage{graphicx}
\usepackage{amsmath}
\usepackage{float}
\setlength{\parindent}{0pt}
\pagestyle{plain}

\begin{document}

% ------------------------------------------------------------------------------
% PORTADA
% ------------------------------------------------------------------------------
\begin{center}
{\Large \textbf{Pontificia Universidad Javeriana}}\\[0.3cm]
Facultad de Ciencias Economicas y Administrativas\\[0.3cm]
Estadistica -- 2026-II\\[1cm]
{\LARGE \textbf{Brecha digital: ?`el acceso a internet y computador afecta el valor agregado universitario?}}\\[0.5cm]
{\large Informe Final del Proyecto}\\[1cm]
\textbf{Grupo 3}\\[0.3cm]
Castillo Lopez Angie Mariana\\
Olaya Esguerra Juan Pablo\\
Novoa Diaz Sebastian\\
Garcia Sanabria Santiago\\[1cm]
27 de noviembre de 2026
\end{center}

\newpage
\tableofcontents
\newpage

% ------------------------------------------------------------------------------
\section{Introduccion}
% ------------------------------------------------------------------------------
% Contexto: brecha digital en Colombia y acceso desigual a tecnologia.
% Pregunta: los estudiantes sin acceso a internet y computador durante
% el bachillerato obtienen menor valor agregado en la universidad?
% Concepto de delta (valor agregado estandarizado) como medida de progreso.

% ------------------------------------------------------------------------------
\section{Datos y Metodologia}
% ------------------------------------------------------------------------------
% Fuente: cruce Saber 11 -- Saber Pro (ICFES).
% Archivos: cruce\_saber11\_saberpro.parquet, saber11.parquet, saber\_pro.parquet.
% Variable dependiente 1: delta = scale(PUNT\_GLOBAL\_sbpro) - scale(PUNT\_GLOBAL\_sb11).
% Variable dependiente 2: MOD\_LECTURA\_CRITICA\_PUNT.
% Predictores: FAMI\_TIENEINTERNET\_sb11, FAMI\_TIENECOMPUTADOR\_sb11,
%   FAMI\_ESTRATOVIVIENDA\_sb11, ESTU\_ZONA\_RESIDE\_sb11 (+ PUNT\_LECTURA\_CRITICA\_sb11 en M2).
% Modelo 1: delta en funcion de acceso digital y contexto.
% Modelo 2: lectura critica SBPro controlando por lectura critica SB11.

% ------------------------------------------------------------------------------
\section{Analisis Exploratorio (Fase 1)}
% ------------------------------------------------------------------------------
% Distribucion de delta por acceso a internet y computador.
% Tablas cruzadas internet/computador vs zona y estrato.
% Estadisticas descriptivas de lectura critica SB11 y SBPro.

% ------------------------------------------------------------------------------
\section{Intervalos de Confianza y Pruebas de Hipotesis (Fase 2)}
% ------------------------------------------------------------------------------
% IC para la media de delta por acceso a internet.
% Prueba t para diferencia de medias delta con vs sin internet.
% Prueba de Wilcoxon como alternativa no parametrica.

% ------------------------------------------------------------------------------
\section{Modelo de Regresion (Fase 3)}
% ------------------------------------------------------------------------------
% Modelo 1: delta $\sim$ internet + computador + estrato + zona.
% Modelo 2: lectura critica SBPro $\sim$ internet + computador + estrato + zona
%   + lectura critica SB11.
% Tabla de coeficientes con IC al 95\%.
% Comparacion de modelos: R2 ajustado, AIC, BIC.
% Predicciones: delta predicho para estudiante con vs sin acceso digital.

% ------------------------------------------------------------------------------
\section{Diagnosticos y Robustez}
% ------------------------------------------------------------------------------
% Diagnosticos para ambos modelos: residuales vs ajustados, QQ-plot.
% Test de Breusch-Pagan (homocedasticidad).
% VIF (multicolinealidad).
% Distancia de Cook (observaciones influyentes).
% Errores robustos HC1: comparacion con errores estandar clasicos.

% ------------------------------------------------------------------------------
\section{Resultados y Discusion}
% ------------------------------------------------------------------------------
% Magnitud del efecto del acceso digital sobre el valor agregado.
% Comparacion de predicciones con vs sin internet.
% Reflexion sobre la brecha digital como barrera a la movilidad educativa.

% ------------------------------------------------------------------------------
\section{Conclusiones}
% ------------------------------------------------------------------------------
% Respuesta a la pregunta de investigacion.
% Limitaciones del estudio.
% Implicaciones para politica de conectividad y acceso digital en educacion.

% ------------------------------------------------------------------------------
\section*{Referencias}
% ------------------------------------------------------------------------------
% ICFES, literatura sobre brecha digital y educacion en Colombia.

% ------------------------------------------------------------------------------
\section*{Anexo: Codigo R}
% ------------------------------------------------------------------------------
% Incluir codigo R utilizado (fase3\_grupo3.R) o referenciarlo.

\end{document}
